\documentclass[11pt]{article}
\usepackage[margin=1in]{geometry}
\usepackage{amsmath,amssymb,bm}
\usepackage{booktabs}
\usepackage{graphicx}
\usepackage{caption}
\usepackage{float}
\usepackage[colorlinks=true,linkcolor=blue,citecolor=blue]{hyperref}
\captionsetup{font=small,labelfont=bf}

\title{Equal-Tailed versus Width-Optimized Conformal Scores\\
under CDF Estimation Error: A Pilot Study}
\author{CKME--CP project --- pilot for the score-choice ablation (E4)}
\date{July 16, 2026 (v3: adds the PCP sub-study; 50 macroreplications throughout)}

\newcommand{\Fhat}{\widehat{F}}

\begin{document}
\maketitle

\begin{abstract}
\noindent
Split conformal prediction guarantees coverage for \emph{any} nonconformity
score; scores differ only in interval width. For skewed output distributions,
a width-optimized score promises intervals up to $12$--$28\%$ narrower than
the equal-tailed default used in our paper. This pilot tests whether that
promise survives when the conditional CDF must be estimated from a realistic
simulation budget. It does not. Across $50$ macroreplications per budget, the
optimizer's \emph{typical} run is only mildly narrower than the equal-tailed
baseline (median $\approx -3\%$, far from the oracle $-12.3\%$), while a
minority of runs fail catastrophically (up to $+224\%$ wider) because the
optimized cut points chase poorly estimated tails. The failures dominate the
mean at both budgets ($+14.0\%$ at $600$ runs, $+7.6\%$ at $2{,}000$ runs);
raising the budget halves the failure rate ($20\% \to 10\%$) but does not
eliminate it. A fixed protective constraint reduces the damage without ever
beating the baseline. A companion sub-study ports the paper's
scale-normalization principle to the sample-based PCP score (Wang et al.,
2023): under heteroscedastic noise, raw PCP tilts conditional coverage
strongly with the noise scale (bin coverage $0.73$--$1.00$, exactly as our
score-homogeneity analysis predicts), and normalizing its score by the
replication-based scale estimate compresses the tilt $4.4$-fold. Both
studies support the paper's central thesis: efficiency and conditional
behavior are governed at the estimation layer, and principles derived there
transfer across score families.
\end{abstract}

\section{Purpose and connection to the paper}

Our paper constructs prediction intervals for stochastic simulation outputs.
A kernel-based estimator (CKME) learns the conditional distribution function
$F(t \mid x)$ of the output $Y$ at input $x$ from a designed set of
simulation runs, and split conformal prediction (CP) calibrates the intervals
on held-out runs so that $90\%$ coverage is guaranteed regardless of how well
$F$ was estimated. The paper's central thesis is a division of labor: the
\emph{calibration layer} guarantees coverage, while interval
\emph{efficiency} (width) is governed by the \emph{estimation layer}.

Within this pipeline, a \emph{conformal score} determines the shape of the
interval. The paper uses the score $\lvert \Fhat(y\mid x) - \tfrac12 \rvert$,
which yields equal-tailed intervals. A natural question, raised by the
advisor, is whether this score is the best choice: for skewed output
distributions, a width-optimized score promises intervals that are
$12$--$28\%$ narrower in theory. This pilot answers two questions
empirically:
\begin{itemize}
  \item[\textbf{Q1}] Do width-optimized score variants realize their
  theoretical gains when $F$ must be estimated from a realistic simulation
  budget?
  \item[\textbf{Q2}] Can a constraint derived from our tail-error analysis
  protect those gains?
\end{itemize}
The results inform the score-choice ablation planned for the paper and
provide a direct test of the layer-separation thesis: promises made at the
score layer can only be cashed at the estimation layer.

\section{Background}

Write $u = \Fhat(y \mid x) \in [0,1]$ for the estimated probability position
of an output value $y$ (the probability integral transform, PIT). If $\Fhat$
were exact, $u$ would be uniformly distributed. A $90\%$ prediction interval
can be built by keeping all $y$ whose PIT falls in a band
$[\theta,\, \theta + 0.9]$, and CP widens or narrows this band by a
calibrated margin $\hat q$ so that $90\%$ coverage holds exactly. Every
choice of the cut point $\theta \in [0, 0.1]$ gives a valid method; the
choices differ only in width:
\begin{itemize}
  \item \textbf{Equal-tailed (ET).} $\theta = 0.05$, cutting $5\%$ from each
  tail; this is the paper's default score
  $\lvert \Fhat(y \mid x) - \tfrac12 \rvert$ in disguise.
  \item \textbf{Width-optimized (opt).} For each $x$, pick $\theta^{*}(x)$
  minimizing the estimated band width
  $\Fhat^{-1}(\theta + 0.9 \mid x) - \Fhat^{-1}(\theta \mid x)$. For the
  skewed test distribution below, the \emph{oracle} version of this rule
  gives intervals $12.3\%$ narrower than ET, by cutting almost everything
  from the long right tail (oracle cut $\theta^{*} \approx 0.003$; see the
  oracle reference widths in Section~3).
  \item \textbf{Bias-constrained opt (b-opt, our proposal).} The same width
  minimization, but interval endpoints may not sit where the estimated
  density is below $15\%$ of its local peak. This is an implementable proxy
  for our theoretical tail-error formula, which shows that quantile
  estimation error is amplified by $1/f$ --- precisely in the deep tails the
  width optimizer likes to chase.
\end{itemize}

The hazard being tested: moving $\theta$ toward $0$ places an endpoint at an
extreme quantile, where the estimated CDF is flattest and least reliable. If
the tail estimate is poor, the calibration layer must widen \emph{all}
intervals to repair coverage, so the optimized score can end up wider than
the simple one.

\section{Experimental design}

\paragraph{Data-generating process.}
$Y = m(x) + 0.3\,\varepsilon$ with $m(x) = e^{x/10}\sin x$ on
$x \in [0, 2\pi]$, and $\varepsilon$ a standardized Gamma$(2)$ variable
(mean $0$, variance $1$, skewness $1.41$). The noise is skewed so that width
optimization has a real prize (oracle: $-12.3\%$ width), and homoscedastic so
that the score comparison is not confounded with the paper's separate
bandwidth-adaptation question.

\paragraph{Oracle reference widths.}
Because the noise is homoscedastic, both oracle widths are constant in $x$
and can be computed exactly from the noise quantile function. Let
$q_{\varepsilon}(p) = \bigl(F^{-1}_{\Gamma(2)}(p) - 2\bigr)/\sqrt{2}$ denote
the quantile function of the standardized Gamma$(2)$ noise. The oracle
\emph{equal-tailed} width at level $1-\alpha = 0.9$ is
\begin{equation}
w^{\mathrm{or}}_{\mathrm{ET}}
  = s \,\bigl[\, q_{\varepsilon}(0.95) - q_{\varepsilon}(0.05) \,\bigr]
  = 0.3 \times 3.103 = 0.931,
\label{eq:oracle-et}
\end{equation}
and the oracle \emph{shortest} width is the best member of the PIT-band
family,
\begin{equation}
w^{\mathrm{or}}_{\mathrm{short}}
  = s \,\min_{\theta \in [0,\,\alpha]}
    \bigl[\, q_{\varepsilon}(\theta + 1 - \alpha) - q_{\varepsilon}(\theta)
    \,\bigr]
  = 0.3 \times 2.721 = 0.816,
\label{eq:oracle-short}
\end{equation}
attained at $\theta^{*} \approx 0.003$ (an interior optimum: the Gamma$(2)$
density also vanishes at the left endpoint of its support, so cutting
exactly $\theta = 0$ is not optimal). Equations~\eqref{eq:oracle-et} and
\eqref{eq:oracle-short} are two views of one comparison:
$w^{\mathrm{or}}_{\mathrm{ET}} / w^{\mathrm{or}}_{\mathrm{short}} - 1 =
+14.0\%$ (ET is $14\%$ wider than the shortest band), equivalently
$1 - w^{\mathrm{or}}_{\mathrm{short}} / w^{\mathrm{or}}_{\mathrm{ET}} =
-12.3\%$ (the shortest band is $12.3\%$ narrower than ET). Since the
experiment reports paired differences \emph{relative to ET}, the relevant
oracle target throughout is $-12.3\%$.

\paragraph{Data sizes.}
Training: $60$ space-filling design sites $\times\, 10$ replications $= 600$
simulator runs (low budget), or $200 \times 10 = 2{,}000$ runs (high budget).
Calibration: $300$ fresh input--output pairs drawn uniformly. Test: $1{,}000$
fresh pairs. The CDF is evaluated on a $220$-point output grid; CKME
hyperparameters are held fixed across arms (length scale $0.3$, ridge
$10^{-3}$, output smoothing bandwidth $0.05$). Macroreplications: $50$ per
budget, each regenerating all data from a fresh seed (seeds $20260716{+}k$
and $20260800{+}k$, $k = 0,\dots,49$). Total compute $\approx 3$ minutes on a
laptop-scale machine.

\paragraph{Protocol.}
Within each macroreplication, all three scores share the \emph{same} fitted
CKME model and the \emph{same} calibration set, so differences between arms
are attributable to the score alone (a paired design). The cut point
$\theta^{*}(x)$ is searched over a $41$-point grid on $[10^{-4}, 0.1]$.
Validity is measured as the coverage of the conformal set itself (the exact
quantity the CP theorem guarantees); width is measured on the
monotone-projected interval, the paper's reporting convention.

\section{Results}

Table~\ref{tab:summary} summarizes all runs; Figures~\ref{fig:validity} and
\ref{fig:paired} show the same information graphically.

\begin{table}[H]
\centering
\small
\caption{Summary over $50$ macroreplications per budget. ``Coverage'' is the
marginal coverage of the conformal set at nominal level $0.90$. ``Width'' is
the mean prediction-interval width in response units; the oracle equal-tailed
width is $0.931$ and the oracle shortest width is $0.816$. ``Paired $\Delta$
width vs.\ ET'' is the per-macroreplication percentage width difference
against the ET arm. ``Fail.\ rate'' is the share of macroreplications more
than $20\%$ wider than ET (catastrophic tail-chasing failures). ``Mean
$\theta^{*}$'' is the average lower-tail cut point chosen by the score (ET
fixes $0.050$; smaller values indicate deeper tail-chasing).}
\label{tab:summary}
\begin{tabular}{llcccccc}
\toprule
Budget & Score & Coverage & Width & Width &
Paired $\Delta$ vs.\ ET & Fail. & Mean \\
 & & & (mean) & (SD) & (mean / median) & rate & $\theta^{*}$ \\
\midrule
$B=600$  & ET    & 0.900 & 0.955 & 0.065 & ---                  & ---    & 0.050 \\
         & opt   & 0.901 & 1.096 & 0.465 & $+14.0\%$ / $-3.0\%$ & $20\%$ & 0.016 \\
         & b-opt & 0.900 & 0.998 & 0.274 & $+4.2\%$ / $-0.7\%$  & $8\%$  & 0.024 \\
\midrule
$B=2000$ & ET    & 0.902 & 0.904 & 0.041 & ---                  & ---    & 0.050 \\
         & opt   & 0.904 & 0.979 & 0.444 & $+7.6\%$ / $-3.1\%$  & $10\%$ & 0.013 \\
         & b-opt & 0.902 & 0.946 & 0.299 & $+4.3\%$ / $-1.4\%$  & $4\%$  & 0.019 \\
\bottomrule
\end{tabular}
\end{table}

\begin{figure}[H]
\centering
\includegraphics[width=\textwidth]{pilot_fig1_validity_width.png}
\caption{Validity and efficiency of the three conformal scores at two
simulation budgets ($50$ macroreplications each). \textbf{(a)}~Marginal
coverage of the conformal prediction set: each colored dot is one
macroreplication, black bars show the mean $\pm$ one standard error, and the
dashed gray line is the nominal level $0.90$. All six arm$\times$budget cells
sit on the nominal line --- score choice never threatens validity.
\textbf{(b)}~Mean interval width (response units): dotted horizontal lines
mark the oracle equal-tailed width (blue, $0.931$) and the oracle shortest
width (red, $0.816$). At both budgets the width-optimized arms show a
compact bulk near the ET cloud plus a scatter of extreme outliers ---
the signature of rare tail-chasing failures.}
\label{fig:validity}
\end{figure}

\begin{figure}[H]
\centering
\includegraphics[width=0.72\textwidth]{pilot_fig2_paired_deltas.png}
\caption{Paired width comparison against the equal-tailed baseline. Each dot
is one macroreplication's percentage width difference
$(\text{arm} - \text{ET})/\text{ET}$; negative values mean the arm is
narrower than ET. Black diamonds: mean $\pm$ one standard error. The red
dotted line ($-12.3\%$) is the oracle gain of the shortest band, i.e., the
best achievable value. The distributions are heavy-tailed upward at
\emph{both} budgets: medians are mildly negative ($-3.0\%$ and $-3.1\%$ for
opt) but $20\%$ ($B=600$) and $10\%$ ($B=2000$) of macroreplications exceed
$+20\%$, with extremes of $+172\%$ and $+224\%$, dragging the means to
$+14.0\%$ and $+7.6\%$. The b-opt constraint reduces the failure rate
($8\%$ and $4\%$) without ever beating ET on average.}
\label{fig:paired}
\end{figure}

\paragraph{Findings.}
\begin{itemize}
  \item[\textbf{F1}] \textbf{(Validity.)} All three scores attain $0.90$
  coverage at both budgets --- as theory predicts, score choice is purely an
  efficiency question.
  \item[\textbf{F2}] \textbf{(The theoretical gain does not survive
  estimation.)} The width optimizer chases the poorly estimated deep tail
  (mean $\theta^{*} = 0.016$ vs.\ $0.050$ for ET). Its risk profile is
  heavy-tailed: the typical macroreplication is mildly narrower than ET
  (median $-3.0\%$, a quarter of the oracle prize), but $20\%$ of
  macroreplications at $B = 600$ come out more than $20\%$ wider (up to
  $+172\%$), inverting the mean to $+14.0\%$ with a width SD seven times
  that of ET.
  \item[\textbf{F3}] \textbf{(More budget shrinks but does not remove the
  failure mode.)} At $B = 2000$ the typical gain is stable (median
  $-3.1\%$; the optimizer is narrower in $70\%$ of macroreplications,
  signed-rank $p = 0.03$), and the failure rate halves ($20\% \to 10\%$) ---
  yet the surviving failures are as extreme as ever (up to $+224\%$), so the
  \emph{mean} remains inverted ($+7.6\%$). A caution on experiment sizing:
  an earlier $8$-macroreplication version of this pilot happened to contain
  no failure at $B = 2000$ and showed a mean gain of $-2.6\%$; with $50$
  macroreplications that conclusion reverses. Rare-failure phenomena require
  large macroreplication counts.
  \item[\textbf{F4}] \textbf{(A fixed protective constraint is not enough.)}
  The density-floor constraint (b-opt) cuts the failure rate at both budgets
  ($20\% \to 8\%$ and $10\% \to 4\%$) and reduces the mean damage, but it
  never beats ET on average ($+4.2\%$ and $+4.3\%$) and its typical gain is
  smaller than the unconstrained optimizer's. A fixed threshold is either
  too loose to prevent failures or too tight to keep the gains; the safe
  tail-chasing depth depends on estimation quality, so the constraint must
  be calibrated from the estimated tail error itself.
\end{itemize}

\section{Extension: a sample-based score (PCP) and scale normalization}

\paragraph{Background and prediction.}
Probabilistic Conformal Prediction (PCP; Wang et al., AISTATS 2023) replaces
CDF evaluation by \emph{sampling}: draw $K$ conditional samples
$\hat y_1,\dots,\hat y_K$ from the estimated conditional law, use the score
$E = \min_k \lvert y - \hat y_k\rvert$, and report the union of $K$ intervals
of calibrated radius. Because our projected CKME CDF can be sampled by
inverse transform, PCP is directly implementable on top of our estimator.
Our score-homogeneity analysis makes a sharp prediction: under the
location-scale model, $E \mid X = x$ is distributed as
$s(x)\cdot\min_k\lvert\varepsilon-\varepsilon_k'\rvert$ --- the score
\emph{scales with the local noise level} --- so a single global threshold
must over-cover where $s(x)$ is small and under-cover where $s(x)$ is large,
the same pathology as fixed-bandwidth DCP. The predicted repair is the
paper's own device: normalize the score by the replication-based scale
estimate, $E/\hat s(x)$ (arm \texttt{pcp\_scaled}).

\paragraph{Design.}
Arms: \texttt{ET} (reference), \texttt{pcp} ($K = 40$ inverse-transform
samples per point), \texttt{pcp\_scaled} ($E/\hat s_{\mathrm{NW}}(x)$, scale
from training-site replications with Nadaraya--Watson smoothing). DGPs: the
skewed Gamma DGP of Section~3 (homoscedastic control) and a heteroscedastic
Gaussian DGP with the same mean function and
$s(x) = 0.10 + 0.20(x-\pi)^2$ (scale ratio $\approx 20$ across the domain).
Budget $B = 600$; calibration $300$; test $2{,}000$ points in $10$
equal-width $x$-bins; $50$ macroreplications per DGP (seeds $20260900{+}k$,
$20261000{+}k$). Set size is the Lebesgue measure of the union (PCP sets are
generally disconnected; Table~\ref{tab:pcp} reports the mean component
count).

\begin{table}[H]
\centering
\small
\caption{PCP sub-study over $50$ macroreplications per DGP. ``Size'' is the
Lebesgue measure of the conformal set (for ET, the interval width).
``\#Comp.''\ is the mean number of disconnected components. ``Bin-cov.\
range'' is $\max - \min$ of the $10$-bin coverage profile and
``$\rho(\mathrm{cov}, s)$'' its Spearman correlation with the noise scale
(heteroscedastic DGP only; the homoscedastic DGP has constant $s$).}
\label{tab:pcp}
\begin{tabular}{llccccc}
\toprule
DGP & Score & Coverage & Size (SD) & \#Comp. & Bin-cov.\ range &
$\rho(\mathrm{cov}, s)$ \\
\midrule
skewed Gamma  & ET          & 0.901 & 0.958 (0.060) & 1.00 & --- & --- \\
(homosced.)   & pcp         & 0.900 & 1.113 (0.051) & 4.48 & --- & --- \\
              & pcp\_scaled & 0.900 & 1.118 (0.052) & 4.32 & --- & --- \\
\midrule
heterosced.\  & ET          & 0.898 & 2.495 (0.143) & 1.00 & 0.109 & $-0.86$ \\
Gaussian      & pcp         & 0.902 & 2.905 (0.163) & 3.13 & \textbf{0.274} & $-0.98$ \\
              & pcp\_scaled & 0.897 & 3.004 (0.135) & 3.96 & \textbf{0.062} & $-0.81$ \\
\bottomrule
\end{tabular}
\end{table}

\begin{figure}[H]
\centering
\includegraphics[width=\textwidth]{pilot_pcp_fig.png}
\caption{PCP sub-study. \textbf{(a)}~Bin coverage of the conformal set on
the heteroscedastic Gaussian DGP: each colored line is one score (markers =
bin means over $50$ macroreplications, error bars = one standard error), the
dashed gray line is the nominal level $0.90$, and the dotted gray curve
(right axis) is the noise scale $s(x)$. Raw PCP (red) mirrors $s(x)$:
over-coverage up to $0.999$ where the noise is small, under-coverage down to
$0.725$ where it is large. Scale normalization (green, $E/\hat s(x)$)
flattens the profile ($4.4\times$ smaller range), flatter even than the
fixed-bandwidth ET reference (blue). \textbf{(b)}~Mean set size
(Lebesgue measure, mean $\pm$ one standard error) by score and DGP.}
\label{fig:pcp}
\end{figure}

\paragraph{Findings.}
\begin{itemize}
  \item[\textbf{F5}] \textbf{(Validity.)} All arms attain marginal coverage
  $0.90$ on both DGPs --- sample-based scores and normalization leave the
  guarantee untouched.
  \item[\textbf{F6}] \textbf{(The tilt prediction is confirmed, and our
  normalization repairs it.)} Raw PCP's bin coverage tracks $-s(x)$ almost
  perfectly (Spearman $-0.98$) with a $27$-percentage-point range;
  $E/\hat s(x)$ compresses the range to $6$ points, flatter than the
  fixed-bandwidth ET reference. This is direct evidence that the paper's
  estimation-layer principle (align the score's implicit scale with the
  local aleatoric scale) transfers to other score families.
  \item[\textbf{F7}] \textbf{(PCP is a baseline, not a replacement.)} On
  these unimodal DGPs PCP pays a $16$--$20\%$ measure premium over the ET
  band and fragments into $3$--$4.5$ components on average; its structural
  advantage (multimodal conditional laws) is absent here. Scale
  normalization slightly increases average measure on the heteroscedastic
  DGP ($3.00$ vs $2.91$) --- reallocating coverage from narrow over-covered
  regions to wide under-covered ones raises total measure, the same
  price/benefit structure as adaptive vs fixed bandwidth.
\end{itemize}

\section{Conclusion}

Score choice never affects the coverage guarantee, but its efficiency
promises are only as good as the estimated CDF behind them. Against an
oracle prize of $-12.3\%$, the width-optimized score delivers a typical gain
of only $\approx 3\%$ and a heavy tail of catastrophic failures (up to
$+224\%$ wider) that inverts its mean performance at both budgets tested;
more data thins the failures out but does not remove them. A fixed
protective constraint reduces the damage without ever winning. The PCP
sub-study shows the same thesis from the constructive side: a sample-based
score inherits a scale-inhomogeneity exactly where our analysis predicts,
and the paper's replication-based normalization repairs it while leaving the
guarantee untouched. Practical default for the paper: keep the equal-tailed
score; report the PCP arms as portability evidence.

\section{Implications for the paper}

\begin{itemize}
  \item Keep the equal-tailed DCP score as the paper's main score; this
  pilot justifies the choice quantitatively rather than by convention.
  \item Report the heavy-tailed risk profile of width optimization
  (Figure~\ref{fig:paired}) as part of the score-choice ablation: it is
  direct evidence for the paper's thesis that score-layer efficiency
  promises are bounded by estimation-layer quality --- the same argument
  that motivates our scale-adaptive bandwidth.
  \item An error-calibrated tail-depth rule (replacing the fixed density
  floor by the paper's tail-error formula, including its variance part) is a
  promising follow-up, but it is future work, not a claim of the current
  paper.
  \item Methodological note for all experiments: rare-failure effects are
  invisible at small macroreplication counts (the $8$-rep version of this
  pilot reached the opposite conclusion at $B=2000$); the paper's standard
  of $\geq 50$ macroreplications is necessary, not cosmetic.
\end{itemize}

\vspace{1em}
\noindent\small\textit{Reproducibility.}
\texttt{exp\_coverage\_mechanism/pilot\_score\_design/} contains the runners
(\texttt{pilot\_score\_design.py}, \texttt{pilot\_pcp.py}), raw
per-macroreplication results (\texttt{pilot\_n50\_b600.csv},
\texttt{pilot\_n50\_b2000.csv}, \texttt{pilot\_pcp\_n50.csv}), and the figure
scripts (\texttt{plot\_pilot\_report\_figs.py}, \texttt{plot\_pcp\_fig.py}).
Seeds and all settings are hard-coded in the runners; the full study reruns
in about $4$ minutes.

\end{document}
